\chapter{Oscillations}\label{ch:osc}

One equation runs this chapter: $\ddot x+\omega_0^{2}x=0$. Springs,
pendulums, RLC circuits (ch15) and molecule vibrations are all its
instances. Add damping and a drive and the same equation explains resonance
--- from tuning forks to the Tacoma bridge.

\section{SHM: the equation and its family}\label{sec:osc-shm}

\begin{equation}
  \ddot x+\omega_0^{2}x=0\;\Rightarrow\;
  x(t)=A\cos(\omega_0t+\phi),\label{eq:shm}
\end{equation}
two constants $A,\phi$ paid by two initial conditions. The frequency
$\omega_0$ is set by the physics; the amplitude is yours. Notice what's
\emph{not} in the solution: $A$ never changes the period --- the signature
of a linear restoring force.

\section{Spring--mass: the canonical case}\label{sec:osc-spring}

$-kx=m\ddot x$ gives $\omega_0=\sqrt{k/m}$ and
$T=2\pi\sqrt{m/k}$: stiffer is faster, heavier is slower. Energy trades
$\tfrac12kx^{2}$ and $\tfrac12mv^{2}$ with
$E=\tfrac12kA^{2}$ --- quadratic in amplitude, so doubling $A$ quadruples
the energy.

\begin{example}[Energy split in a spring oscillator]\label{ex:springE}
A $0.5$\,kg mass on a $200$\,N/m spring oscillates with $A=0.15$\,m.
Find $E$, $v_{\max}$, and the speed at $x=A/2$.
\par\smallskip
\textbf{Strategy.} $E$ is set by amplitude; at any $x$ the remainder
after $\tfrac12kx^{2}$ is kinetic.
\begin{solution}
$E=\tfrac12(200)(0.0225)=2.25$\,J;
$v_{\max}=\sqrt{2E/m}=\sqrt{9}=3$\,m/s at centre.
At $x=0.075$: $U=\tfrac12(200)(0.0056)=0.56$\,J, so
$v=\sqrt{2(2.25-0.56)/0.5}=2.6$\,m/s --- three quarters of max at half
amplitude: energy spends most of the cycle near the turning points.
\checkmark
\end{solution}
\end{example}

\section{Pendulums: small-angle as an approximation with a limit}\label{sec:osc-pend}

For a pendulum the restoring torque is $-mgL\sin\theta$
(figure~\ref{fig:pend}); linearising $\sin\theta\approx\theta$ gives
$\omega_0=\sqrt{g/L}$ --- good to $1\%$ for
$\theta_0\lesssim23^{\circ}$, and the first correction
$T\approx T_0(1+\theta_0^{2}/16)$ is the hard exercise at the end.
A physical pendulum generalises it:
$\omega_0=\sqrt{mgd/I}$.

\begin{figure}[htbp]\centering
\begin{tikzpicture}[scale=1.2,force/.style={->,thick,accent}]
\fill (0,0) circle (0.04) node[above,font=\scriptsize]{pivot};
\draw[thick] (0,0) -- (1.8,-2.5);
\fill[fill=gray!30] (1.8,-2.5) circle (0.12);
\draw[dashed] (0,0) -- (0,-2.7);
\draw[force] (1.8,-2.5) -- (1.8,-3.4) node[below,font=\scriptsize]{$mg$};
\draw[force,dashed,amber] (1.8,-2.5) -- (1.1,-2.95) node[left,font=\scriptsize]{$mg\sin\theta$};
\draw (0.55,0) arc (0:-54:0.55);
\node[font=\scriptsize] at (0.85,-0.45) {$\theta$};
\node[font=\scriptsize,align=left] at (2.6,-1.5) {restore $\approx-mg\theta$\\for $\theta\lesssim15^{\circ}$};
\node[font=\scriptsize,align=left] at (2.6,-1.5) {};
\end{tikzpicture}

\caption{The pendulum's restoring force $-mg\sin\theta$ is \emph{almost}
linear --- small-angle is an approximation with a budget, not a property.}
\label{fig:pend}
\end{figure}

\begin{example}[Meter-stick physical pendulum]\label{ex:physpend}
A uniform meter stick pivots at one end. Find its period.
\par\smallskip
\textbf{Strategy.} Physical pendulum: $\omega_0^{2}=mgd/I$ with $d=L/2$,
$I=\tfrac13mL^{2}$.
\begin{solution}
$\omega_0^{2}=\frac{g(L/2)}{L^{2}/3}=\frac{3g}{2L}$, so
$T=2\pi\sqrt{\frac{2L}{3g}}=2\pi\sqrt{\frac{2}{29.4}}=1.64$\,s. A simple
pendulum this long would take $2.0$\,s --- the stick's spread mass makes
it faster. \checkmark
\end{solution}
\end{example}

\section{Damped oscillations: three regimes}\label{sec:osc-damped}

Add $-b\dot x$: $\ddot x+\gamma\dot x+\omega_0^{2}x=0$ with
$\gamma=b/m$. The discriminant sorts the behaviour
(figure~\ref{fig:damp}): \textbf{underdamped} ($\gamma<2\omega_0$)
oscillates inside an $e^{-\gamma t/2}$ envelope at
$\omega_d=\sqrt{\omega_0^{2}-\gamma^{2}/4}$; \textbf{critical}
($\gamma=2\omega_0$) returns fastest without ringing;
\textbf{overdamped} ($\gamma>2\omega_0$) creeps back slowly. Car
suspension and door closers are engineered to sit just under critical.

\begin{figure}[htbp]\centering
\begin{tikzpicture}
\begin{axis}[width=9cm,height=5cm,xlabel={$t$},ylabel={$x$},
xmin=0,xmax=20,ymin=-1.3,ymax=1.3,xtick={0,5,10,15,20},font=\small,legend pos=north east]
\addplot[very thick,accent,domain=0:20,samples=120] {exp(-0.08*x)*cos(deg(2*x))};
\addlegendentry{under}
\addplot[very thick,amber,domain=0:20,samples=60] {exp(-0.5*x)*cos(deg(1.6*x))};
\addlegendentry{critical-ish}
\addplot[very thick,dashed,domain=0:20,samples=60] {exp(-0.9*x)};
\addlegendentry{over}
\end{axis}
\end{tikzpicture}

\caption{The three damping regimes. Notice only the underdamped case
crosses zero --- the others surrender exponentially.}
\label{fig:damp}
\end{figure}

\begin{example}[Extracting the constants from a trace]\label{ex:damped}
A decaying trace shows maxima at $1.00$, $0.82$, $0.67$ at $t=0$, $2.1$,
$4.2$\,s. Find $\gamma$, $\omega_d$, $\omega_0$.
\par\smallskip
\textbf{Strategy.} Envelope ratio per period gives $\gamma$; the period
gives $\omega_d$; the relation $\omega_d^{2}=\omega_0^{2}-\gamma^{2}/4$
closes it.
\begin{solution}
Period $T_d\approx2.1$\,s $\Rightarrow\omega_d=2.99$\,rad/s.
Ratio per period $0.82$: $e^{-\gamma T_d/2}=0.82 \Rightarrow
\gamma=0.19$\,s$^{-1}$.
$\omega_0=\sqrt{\omega_d^{2}+\gamma^{2}/4}=\sqrt{8.94+0.009}=2.99$\,rad/s
--- damping barely shifts the frequency here. \checkmark
\end{solution}
\end{example}

\begin{exercise}[easy]
Show the same trace's next two maxima and estimate them. (AP-C)
\end{exercise}

\section{Driven oscillations and resonance}\label{sec:osc-driven}

A drive $F_0\cos\omega_dt$ produces steady amplitude
\begin{equation}
  A(\omega_d)=\frac{F_0/m}{\sqrt{(\omega_0^{2}-\omega_d^{2})^{2}
  +(\gamma\omega_d)^{2}}},\label{eq:amp}
\end{equation}
peaking near $\omega_d=\sqrt{\omega_0^{2}-\gamma^{2}/2}$
(figure~\ref{fig:resonance}). The peak height is set by the
\emph{quality factor} $Q=\omega_0/\gamma$: high-$Q$ systems ring loudly
and narrowly --- the tuning-fork pair of figure~\ref{fig:fork} and, minus
the luck, the Tacoma Narrows bridge.

\begin{example}[Off-resonance driving]\label{ex:offres}
The oscillator of the code run ($\omega_0=2$, $\gamma=0.4$, $F_0/m=1$)
is driven at $\omega_d=3$. Find its steady amplitude and compare with
the peak.
\par\smallskip
\textbf{Strategy.} Plug into \eqref{eq:amp} twice; the ratio is the
point.
\begin{solution}
$A(3)=1/\sqrt{26.4}=0.195$ (from
$\sqrt{(4-9)^{2}+(0.4\cdot3)^{2}}$);
$A(1.98)=1/(0.4\times1.98)=1.26$.
The peak is $6.5\times$ louder --- a $1.5\times$ frequency miss costs
$85\%$ of the response. \checkmark
\end{solution}
\end{example}

\begin{figure}[htbp]\centering
\begin{tikzpicture}
\begin{axis}[width=8.8cm,height=5.5cm,xlabel={$\omega_d$ (rad/s)},ylabel={$A(\omega_d)$},
xmin=0.5,xmax=3.5,ymin=0,ymax=1.6,font=\small,legend pos=north east]
\addplot[very thick,accent,domain=0.5:3.5,samples=120] {1/sqrt((4-x^2)^2+(0.4*x)^2)};
\addlegendentry{$\gamma=0.4$}
\addplot[very thick,amber,domain=0.5:3.5,samples=120] {1/sqrt((4-x^2)^2+(1.2*x)^2)};
\addlegendentry{$\gamma=1.2$}
\addplot[dashed] coordinates {(1.98,0) (1.98,1.6)};
\node[font=\scriptsize] at (axis cs:2.4,1.35) {peak near $\sqrt{\omega_0^2-\gamma^2/2}$};
\end{axis}
\end{tikzpicture}

\caption{Amplitude vs drive frequency for two dampings. Lower $\gamma$
raises and sharpens the peak --- the $Q$ trade-off.}
\label{fig:resonance}
\end{figure}

\begin{example}[Bandwidth from a resonance curve]\label{ex:qfactor}
A resonance curve peaks at $195$\,Hz with amplitude $A_0$ and falls to
$A_0/\sqrt{2}$ at $190$ and $200$\,Hz. Find $Q$ and $\gamma$.
\par\smallskip
\textbf{Strategy.} Half-power bandwidth $\Delta\omega=\gamma$; then
$Q=\omega_0/\gamma$.
\begin{solution}
$\Delta f=10$\,Hz $\Rightarrow \gamma=2\pi\times10=62.8$\,s$^{-1}$,
$Q=195/10=19.5$. Narrow and loud. \checkmark
\end{solution}
\end{example}

\begin{lstlisting}[language=Python,caption={\texttt{code/ch07\_resonance\_sweep.py} --- RK4 driven oscillator, amplitude vs theory.},label={lst:sweep}]
def acc(x_, v_, t_):
    return F0M*math.cos(wd*t_) - GAMMA*v_ - W0*W0*x_
\end{lstlisting}

Output of a real run:
\begin{lstlisting}[numbers=none]
gamma=0.4, w0=2.0: theory peak at sqrt(w0^2-g^2/2)=1.980 rad/s
   wd   A_meas  A_theory    err
  1.2    0.384     0.384   0.0%
  1.6    0.635     0.635   0.0%
  1.9    1.171     1.171   0.0%
  2.0    1.250     1.250   0.0%
  2.1    1.070     1.070   0.0%
  2.4    0.499     0.499   0.0%
  3.0    0.194     0.194   0.0%
Q = w0/gamma = 5.0; A(w_res)/A(0) = 5.03 (theory)
\end{lstlisting}

RK4 tracks the analytic curve to $0.0\%$ at every sampled $\omega_d$ ---
the integrator is now trustworthy enough to trust figure~\ref{fig:resonance}.
Figure~\ref{fig:amplresp} overlays the measured amplitudes on the theory
curve: the agreement is exact, and the peak sits at
$\omega_r=\sqrt{\omega_0^2-\gamma^2/2}=1.980$ --- just below $\omega_0$,
the driven-oscillator subtlety that trips up most first passes.

\begin{figure}[htbp]\centering
\begin{tikzpicture}
\begin{axis}[width=8.8cm,height=5.2cm,xlabel={$\omega_d$ (rad/s)},ylabel={$A(\omega_d)$},
xmin=1,xmax=3.1,ymin=0,ymax=1.45,font=\small,legend pos=north west]
\addplot[very thick,accent,domain=1:3.1,samples=120] {1/sqrt((4-x^2)^2+0.16*x^2)};
\addlegendentry{theory, $\gamma=0.4$, $\omega_0=2$}
\addplot[only marks,mark=*,mark size=2.4pt,amber] coordinates
{(1.2,0.384) (1.6,0.635) (1.9,1.171) (2.0,1.250) (2.1,1.070) (2.4,0.499) (3.0,0.194)};
\addlegendentry{RK4 measured}
\addplot[dashed] coordinates {(1.98,0) (1.98,1.45)};
\node[font=\scriptsize,anchor=west] at (axis cs:1.82,0.15) {$\omega_r=1.980$};
\end{axis}
\end{tikzpicture}

\caption{Simulated steady-state amplitudes (markers) on the analytic
response curve for $\gamma=0.4$, $\omega_0=2$. The maximum is at
$\omega_r=1.980$, not at $\omega_0$ --- damping pulls the peak left.}
\label{fig:amplresp}
\end{figure}

\begin{figure}[htbp]\centering
\begin{tikzpicture}[scale=1.0,font=\scriptsize]
\draw[very thick] (0,0) -- (0,2.4);
\draw[very thick] (0.6,0) -- (0.6,2.4);
\draw[very thick] (0,0) -- (0.6,0);
\draw[dashed,->] (-0.3,2.2) arc (90:200:0.5);
\draw[dashed,->] (0.9,2.2) arc (90:-20:0.5);
\node at (0.3,-0.5) {tuning fork};
\draw[->,thick,amber] (1.6,1.2) -- (3.0,1.2) node[right]{$\omega_d=\omega_0$};
\draw[very thick] (3.6,0) -- (3.6,1.8);
\draw[very thick] (4.2,0) -- (4.2,1.8);
\draw[very thick] (3.6,0) -- (4.2,0);
\draw[dashed] (3.3,1.6) arc (90:210:0.35);
\draw[dashed] (4.5,1.6) arc (90:-30:0.35);
\node at (3.9,-0.5) {same note $\Rightarrow$ rings};
\node at (2.1,-1.4) {sympathetic resonance: only the matching fork answers};
\end{tikzpicture}

\caption{Sympathetic resonance: only the fork sharing $\omega_0$ answers.
Every band's instrument-tuning trick, and every bridge's nightmare.}
\label{fig:fork}
\end{figure}

\section{Phase space: the state is a point}\label{sec:osc-phase}

Plotting $\dot x$ vs $x$ (figure~\ref{fig:phase}) turns the oscillator's
history into a curve: the undamped case draws a closed ellipse, damping
spirals it to the origin. Nonlinear oscillators live here too ---
chapter~17's pendulum without the small-angle crutch is a phase portrait
before it's a formula.

\begin{figure}[htbp]\centering
\begin{tikzpicture}[scale=0.9]
\draw[->] (-3.2,0) -- (3.4,0) node[right,font=\small]{$x$};
\draw[->] (0,-2.4) -- (0,2.6) node[above,font=\small]{$\dot x$};
\draw[very thick,accent] (2.8,0) ellipse (0.001 and 0.001);
\foreach \s in {2.8,2.2,1.6,1.0,0.5}{\draw[accent!60] (0,0) ellipse ({\s} and {\s*0.72});}
\draw[->,thick,amber] (2.6,-1.5) .. controls (1.6,-1.9) and (0.9,-1.1) .. (0.35,-0.35);
\draw[->,thick,amber] (-2.4,1.3) .. controls (-1.5,1.8) and (-0.8,1.0) .. (-0.3,0.3);
\node[font=\scriptsize] at (2.0,2.2) {spiral in: $E$ leaves};
\draw[very thick,dashed] (0,0) ellipse (2.5 and 1.8);
\node[font=\scriptsize] at (-1.9,-2.0) {dashed: undamped ellipse};
\end{tikzpicture}

\caption{Phase portrait: the undamped ellipse versus the damped spiral.
The state's whole future is visible in one picture.}
\label{fig:phase}
\end{figure}

\begin{lstlisting}[language=Python,caption={\texttt{code/ch07\_driven\_transient.py} --- driven oscillator: the beat ring-up dies out, leaving the theory's amplitude.},label={lst:transient}]
m, k, b, F0, w = 1.0, 100.0, 1.0, 1.0, 8.0
w0 = np.sqrt(k / m)
A_theory = F0 / np.sqrt((m * (w0**2 - w**2))**2 + (b * w)**2)
dt, n = 1e-3, 60000
x, v, peaks = 0.0, 0.0, []
for i in range(1, n):
    a = (F0 * np.cos(w * i * dt) - k * x - b * v) / m
    v += a * dt
    x += v * dt
    if i > 50000 and i % 10 == 0:
        peaks.append(x)
\end{lstlisting}

Output of a real run:
\begin{lstlisting}[numbers=none]
w0 = 10.0 rad/s, drive w = 8.0 rad/s
theory steady amplitude: 0.027 m
late-time window max |x|: 0.027 m
beat envelope frequency |w0-w|/2 = 1.0 rad/s
\end{lstlisting}

Off-resonance driving gives a transient beat --- natural frequency
plus drive frequency interfering --- that damping erases in roughly
$\tau=2m/b$; what survives is the steady-state amplitude the formula
promised, to all three printed digits. The simulation's beat
envelope sits at $|\omega_0-\omega|/2=1$\,rad/s, the same
superposition arithmetic chapter~16 puts in your ear.

\section{In engineering practice: vibration budgets}\label{sec:osc-eng}

Every engineering discipline owns a version of equation~\eqref{eq:amp}:
buildings get tuned mass dampers (a secondary oscillator eating the
primary's peak), hard drives sit on elastomer mounts that push
$\omega_0$ below the drive frequencies, and shock absorbers are the
critical-damping spec of \S\ref{sec:osc-damped} sold by the pair. The
common move is always the same: place $\omega_0$ far from the
disturbance's $\omega_d$, then buy the remaining safety factor with
damping --- the two knobs this chapter names.

\section{Common mistakes}\label{sec:osc-mistakes}

\begin{tripwire}[Small-angle used where $\theta$ isn't]
The linearisation costs $\sim1\%$ per $20^{\circ}$ beyond small; at
$60^{\circ}$ the period is $7\%$ long and the whole section's formulas
need the finite-amplitude correction.
\end{tripwire}
\begin{tripwire}[$\omega$ vs $f$]
$\omega=2\pi f$: $T=2\pi\sqrt{m/k}$ is a period, $\sqrt{k/m}$ is a
frequency \emph{in radians}. Mixing them is a factor-$2\pi$ error every
time.
\end{tripwire}
\begin{tripwire}[Amplitude vs energy]
$E=\tfrac12kA^{2}$ is quadratic: ``twice the amplitude'' is four times
the energy, and damping's decay rate is set on the amplitude by the
envelope, not the peak.
\end{tripwire}
\begin{tripwire}[Driving at $\omega_0$ vs $\omega_d$]
Resonance peaks \emph{below} $\omega_0$ when damped ---
$\sqrt{\omega_0^{2}-\gamma^{2}/2}$ --- and the response lags the drive by
$\pi/2$ at the peak.
\end{tripwire}
\begin{tripwire}[``Damping slows the period'' assumed or denied]
It does --- $\omega_d<\omega_0$ --- but weakly ($\gamma^{2}/4$ term), so
check the size before reporting either way.
\end{tripwire}

\section{Summary}\label{sec:osc-summary}

\begin{table}[ht]\centering
\begin{tabularx}{\linewidth}{@{}lY@{}}
\toprule
\tablehead\thead{Idea} & \thead{One-line version}\\
\midrule
SHM & $\ddot x+\omega_0^{2}x=0 \Rightarrow x=A\cos(\omega_0t+\phi)$ \\
spring & $\omega_0=\sqrt{k/m}$; $E=\tfrac12kA^{2}$ \\
pendulum & $\omega_0=\sqrt{g/L}$ small-angle; physical $\sqrt{mgd/I}$ \\
damping & under/critical/over by $\gamma$ vs $2\omega_0$ \\
driven & peak at $\sqrt{\omega_0^{2}-\gamma^{2}/2}$; $Q=\omega_0/\gamma$ \\
phase space & ellipse undamped, spiral damped \\
\bottomrule
\end{tabularx}
\caption{Oscillations cheat sheet.}
\label{tab:oscsummary}
\end{table}

\begin{wins}[Definition of done]
You can now: write and solve the SHM equation; scale spring and pendulum
periods; classify damping regimes and extract constants from a decay
trace; compute resonance peaks and read $Q$ off a bandwidth; and draw the
phase portrait.
\end{wins}

\chapternotes{%
The oscillator-first treatment follows MIT 8.01SC's oscillation module
\cite{ocw801} and AP Physics C unit 7 \cite{apced}; Feynman vol.~I
\S21--24 \cite{feynman1} is the classic resonance narrative. Morin ch.~4
\cite{morin} is the source of the physical-pendulum and finite-amplitude
exercises. The sweep numbers are the listing's real output.}

\begin{rulebox}[{Watch instead}]
\begin{enumerate}
\item \emph{Hooke's Law, Springs, Pendulums, Simple Harmonic Motion} ---
Walter Lewin, MIT 8.01x Lecture 10, 47:41,
\url{https://youtu.be/tNpuTx7UQbw} ---
\S\ref{sec:osc-shm}--\S\ref{sec:osc-pend} at full length (accessed
2026-10-11).
\item \emph{SHM of spring-mass oscillators} --- Khan Academy, 15:13,
\url{https://youtu.be/eRZSzGKDfwU} --- \S\ref{sec:osc-spring} worked
quickly (accessed 2026-10-11).
\end{enumerate}
\end{rulebox}

\section*{Further reading}
\addcontentsline{toc}{section}{Further reading}

\begin{itemize}
\item MIT OpenCourseWare, \emph{8.01SC Classical Mechanics}, oscillations
module --- SHM, damping and resonance with recitations. At
\url{https://ocw.mit.edu}, course ``8.01SC'' (verified 2026-10-11).
\item OpenStax, \emph{College Physics 2e}, ch.~15 ``Oscillatory Motion''
--- free text at \url{https://openstax.org} (verified 2026-10-11).
\item College Board, \emph{AP Physics C: Mechanics CED}, Unit 7 --- the
exam blueprint for SHM. At \url{https://apstudents.collegeboard.org}
(verified 2026-10-11).
\item R.~P. Feynman, \emph{Lectures on Physics}, vol.~I, \S21--24 ---
resonance as narrative; free at Caltech's site.
\item D.~Morin, \emph{Introduction to Classical Mechanics}, ch.~4 ---
oscillations at the harder level; ISBN 978-0521876223.
\item Khan Academy, \emph{Simple harmonic motion} --- drills matching
\S\ref{sec:osc-shm}--\S\ref{sec:osc-driven}, at
\url{https://www.khanacademy.org} under Science $\to$ Physics (verified
2026-10-11).
\end{itemize}

\section*{Additional exercises}
\addcontentsline{toc}{section}{Additional exercises}

\begin{exercise}[easy]
A $0.5$\,kg mass on a $200$\,N/m spring: compute $\omega_0$, $f$, $T$.
(AP-C)
\end{exercise}
\begin{exercise}[easy]
Doubling a spring oscillator's amplitude multiplies its energy by what?
Its maximum speed by what? (AP-C)
\end{exercise}
\begin{exercise}[easy]
A pendulum's $T$ is $2.0$\,s on Earth. Find it on the Moon
($g=1.6$\,m/s$^{2}$). (F=ma)
\end{exercise}
\begin{exercise}[easy]
Convert: an oscillator at $440$\,Hz --- find $\omega$ and $T$. (AP-C)
\end{exercise}
\begin{exercise}[easy]
A $1$\,kg mass on a spring oscillates with $A=0.2$\,m, $T=1$\,s. Find
$k$ and $v_{\max}$. (AP-C)
\end{exercise}
\begin{exercise}[medium]
Two springs in series vs parallel: show effective $k$ is
$k_1k_2/(k_1+k_2)$ vs $k_1+k_2$, and find each for $k_1=100$,
$k_2=300$\,N/m. (F=ma)
\end{exercise}
\begin{exercise}[medium]
A rod of length $L$ pivots at a point $L/4$ from one end (physical
pendulum). Find $T$ and compare with the end-pivoted case. (AP-C FR-style)
\end{exercise}
\begin{exercise}[medium]
At what phase of its cycle does an oscillator carry equal $K$ and $U$?
Give the $x$ value in terms of $A$. (AP-C)
\end{exercise}
\begin{exercise}[medium]
Decide the regime: $m=1$, $k=100$, (a) $b=5$, (b) $b=20$, (c) $b=40$
(all SI) --- and find $\omega_d$ where it exists. (AP-C)
\end{exercise}
\begin{exercise}[medium]
An oscillator loses $20\%$ of its amplitude per period. Find the energy
loss per period and $\gamma$. (F=ma)
\end{exercise}
\begin{exercise}[medium]
A driven oscillator has $\omega_0=10$, $\gamma=2$ (SI). Find the
resonant $\omega_d$, the peak amplitude for $F_0/m=5$, and $Q$. (AP-C)
\end{exercise}
\begin{exercise}[hard]
Use $T\approx T_0(1+\theta_0^{2}/16)$: find the amplitude where a
grandfather clock's period is off by $0.5\%$, and the amplitude for a
$10\%$ slowdown. (USAPhO-style)
\end{exercise}
\begin{exercise}[hard]
A particle sits in $U(x)=-V_0\cos(kx)$. Linearise about a minimum to
get $\omega_0$, then about a maximum to confirm instability.
(APhO-style)
\end{exercise}
\begin{exercise}[medium]
Two oscillators: $x_1=A\cos\omega t$, $x_2=A\cos(\omega t+\pi/2)$ ---
sketch both in phase space and say what the $\pi/2$ buys. (F=ma)
\end{exercise}
\begin{exercise}[medium]
A car's suspension should return to level fastest without ringing:
given $m=300$\,kg and $k=30{,}000$\,N/m, find the critical $b$.
(engineering)
\end{exercise}
